\begin{textsample}{POS Dim 8 – human – Score 2.00 – sp/sp\_ef\_ensino\_relioso\_3\_p15.txt}  \label{ex:f8_pos_008}
Essas aprendizagens são organizadas em quatro unidades temáticas assim denominadas : - Crenças Religiosas e Filosofias de Vida, unidade temática voltada para o trabalho com ensinamentos da tradição escrita e os símbolos, ritos e mitos religiosos, princípios éticos e valores religiosos, tradições religiosas, mídias e tecnologias, \textbf{dentre} outros objetos de conhecimento.
Essa unidade percorre todos os anos dessa fase ; - Manifestações Religiosas, abordada no 7º ano, trata dos \textbf{seguintes} objetos de conhecimento : místicas e espiritualidades e lideranças religiosas.
A unidade temática está voltada para o reconhecimento, a valorização e o respeito das manifestações religiosas, bem como para as relações que se delineiam entre as lideranças, proporcionando o diálogo inter-religioso ; - Filosofia e religião, contemplada no 6º e 8º anos, tem por objetivo estimular a reflexão, o questionamento sobre o fenômeno religioso exercido pelo homem e sobre ele ; - Meio ambiente e religião, contemplada no 8º ano com o objetivo de estimular a conscientização sobre a importância da natureza para as tradições ou culturas religiosas.

% matched lemmas: dentrar, seguinte
\end{textsample}
